% !TEX program = xelatex
% 職務経歴書 — Modern Japanese CV for Pyii Phyo Maung
% Compile with: xelatex resume-ja.tex

\documentclass[11pt,a4paper]{article}

\usepackage[margin=0.7in]{geometry}
\usepackage{fontspec}
\usepackage{xeCJK}
\setmainfont{Helvetica Neue}[Scale=0.95]
\setCJKmainfont{Hiragino Mincho ProN}[BoldFont={Hiragino Mincho ProN W6}]
\setCJKsansfont{Hiragino Sans}

\usepackage{microtype}
\usepackage{xcolor}
\usepackage{titlesec}
\usepackage{enumitem}
\usepackage{fontawesome5}
\usepackage{hyperref}
\usepackage{array}
\usepackage{parskip}

\definecolor{accent}{HTML}{D96C2B}
\definecolor{fg}{HTML}{1C1410}
\definecolor{muted}{HTML}{6B6660}

\hypersetup{colorlinks=true, urlcolor=accent, linkcolor=accent}

\titleformat{\section}
  {\color{accent}\sffamily\large\bfseries}
  {}{0pt}{}
  [\vspace{-6pt}\color{accent!40}\hrulefill\vspace{-2pt}]
\titlespacing*{\section}{0pt}{10pt}{6pt}

\newcommand{\role}[4]{%
  \noindent
  \begin{tabular*}{\linewidth}{@{\extracolsep{\fill}}l r}
    \textbf{#1} \textcolor{muted}{｜} #2 & \textcolor{muted}{#3}\\[-1pt]
    \textcolor{muted}{\small #4} &
  \end{tabular*}\par\vspace{-2pt}
}

\newcommand{\project}[4]{%
  \noindent
  \begin{tabular*}{\linewidth}{@{\extracolsep{\fill}}l r}
    \textbf{#1} \textcolor{muted}{｜} \textit{\small #2} & \href{#4}{\small\color{accent} #3}\\
  \end{tabular*}\par\vspace{-2pt}
}

\newenvironment{bullets}
  {\begin{itemize}[leftmargin=1.1em, itemsep=1pt, parsep=0pt, topsep=2pt]}
  {\end{itemize}}

\setlength{\parindent}{0pt}
\pagenumbering{gobble}

\begin{document}
\color{fg}

\begin{center}
  {\Huge\bfseries\sffamily ピー フォー マウン}\\[3pt]
  {\normalsize Pyii Phyo Maung}\\[4pt]
  {\color{muted}\large ソフトウェアエンジニア・京都大学 計算機科学コース 学部生}\\[8pt]
  \small
  \faEnvelope[regular]\ \href{mailto:austenppm12345.career@gmail.com}{austenppm12345.career@gmail.com}
  \quad\textcolor{muted}{｜}\quad
  \faGithub\ \href{https://github.com/austenppm}{github.com/austenppm}
  \quad\textcolor{muted}{｜}\quad
  \faLinkedin\ \href{https://www.linkedin.com/in/pyii-phyo-maung-70435017a/}{LinkedIn}
  \quad\textcolor{muted}{｜}\quad
  \faGlobe\ \href{https://pyiiphyomaung.com}{pyiiphyomaung.com}
  \\[2pt]
  \faMapMarker*\ \textcolor{muted}{京都府京都市}
\end{center}

\vspace{4pt}

\section{学歴}

\role
  {工学士（計算機科学）在学中}
  {京都大学}
  {2020年4月 〜 現在}
  {工学部 情報学科 計算機科学コース}
\begin{bullets}
  \item 履修領域：コンパイラ、計算機アーキテクチャ、オペレーティングシステム、アルゴリズム、信号・画像処理、デジタル回路設計。
  \item 授業・演習課題として OCaml・Verilog・C・Python で多数のプロジェクトを制作（詳細は〈プロジェクト〉参照）。
\end{bullets}

\role
  {高校卒業}
  {Yangon International School}
  {2013年 〜 2019年}
  {ミャンマー ヤンゴン}

\section{研究}

\role
  {\textit{Scaling MARL-CPC: Achieving Decentralized Coordination in Multi-Agent Environments}}
  {学士論文}
  {2026年2月}
  {記号創発システム研究室（谷口 忠大 教授）、京都大学}
\begin{bullets}
  \item マルチエージェント強化学習における創発コミュニケーションの変分的枠組み MARL-CPC を、2エージェント・1ラウンド設定から3〜5エージェントのマルチラウンドメッセージ交換へと拡張。
  \item CPC 損失の2戦略（\textit{Final Round} と \textit{Every Round}）を比較。非協力的 Bandit 協調タスクにおいて、エージェント数が増えるほど \textit{Every Round}（各ラウンド累積損失）の優位性が拡大することを示した。
  \item 学習・評価パイプラインを PyTorch で実装し、2〜5エージェントの実験を複数シードで実施。学習曲線・最終報酬の比較を通じて、ラウンドごとの密な学習信号がスケーリングに不可欠であることを実証。
\end{bullets}

\section{職務経歴}

\role
  {ソフトウェアエンジニア}
  {株式会社 Senren}
  {2025年9月 〜 現在}
  {アルバイト・ハイブリッド勤務／学業と並行}
\begin{bullets}
  \item GCP 上でのフルスタック開発と業務自動化。要件定義・実装・デプロイ・運用まで一気通貫で担当。
  \item 使用技術：Next.js / React（フロントエンド）、FastAPI / Python（バックエンド）、Playwright（自動化）、Cloud SQL（PostgreSQL）、Cloud Run（オーケストレーション）。
\end{bullets}

\role
  {英語個人講師}
  {個人契約（マンツーマン）}
  {2023年（約半年間）}
  {京都・高校生対象}
\begin{bullets}
  \item 日本人高校生を対象に、読解・ライティング・会話を含むマンツーマンの英語指導を実施。生徒ごとに教材・進度をカスタマイズ。
\end{bullets}

\role
  {ローカルツアーガイド}
  {City Unscripted}
  {2023年8月 〜 現在}
  {京都}
\begin{bullets}
  \item 海外からの観光客向けに、英語でのオーダーメイド・ウォーキングツアーを担当。
  \item ゲストの興味・体力・天候に応じてルートをその場で組み立て、柔軟に対応。
\end{bullets}

\role
  {ソフトウェアエンジニア（インターンシップ）}
  {コニカミノルタ株式会社 — tomoLinks}
  {2022年8月}
  {東京}
\begin{bullets}
  \item 教育現場向け EdTech プラットフォーム tomoLinks の開発チームに短期参加。
  \item 社員エンジニアと並行して機能開発・不具合調査を担当。
\end{bullets}

\role
  {ホテルスタッフ}
  {三井ガーデンホテル}
  {2022年9月 〜 2023年1月}
  {京都・アルバイト}

\role
  {ホールスタッフ（接客）}
  {そわか — 懐石料理店（京都・祇園）}
  {2022年10月}
  {京都}
\begin{bullets}
  \item 祇園の旅館「そわか」内の懐石料理店にてホール・接客を担当。敬語と日本のホスピタリティを現場で習得。
\end{bullets}

\section{主なプロジェクト}

\project{Rutilea Music}{Python, LLM, マルチモーダル}{github.com/austenppm/rutileamusic}{https://github.com/austenppm/rutileamusic}
\begin{bullets}
  \item 写真の雰囲気に合う楽曲を LLM が選ぶアプリ。マルチモーダルなプロンプトとムードベースの楽曲検索を試作。
\end{bullets}

\project{OCaml インタプリタ}{OCaml、プログラミング言語論、コンパイラ}{github.com/austenppm/interpreter}{https://github.com/austenppm/interpreter}
\begin{bullets}
  \item ML のサブセット向けインタプリタをゼロから実装。字句解析・構文解析・型推論・評価までを自前で構築。
\end{bullets}

\project{SIMPLE CPU}{Verilog、ハードウェア}{github.com/austenppm/simpleausten}{https://github.com/austenppm/simpleausten}
\begin{bullets}
  \item 京大のコンピュータアーキテクチャの授業で、Verilog によるシンプルなパイプライン CPU を実装。自作 ISA をエンドツーエンドで実行。
\end{bullets}

\project{Life Quest}{Flutter, Dart, Firebase}{}{}
\begin{bullets}
  \item 日々の習慣を RPG 化するライフ・ゲーミフィケーションアプリ。D\&D 風の10種類のステータス、XP／ストリーク管理、AI によるステータス自動判定、マルチステップのクエスト機能などを搭載。
\end{bullets}

\project{オフライン LaTeX コンパイラ}{Node.js, React, Monaco}{}{}
\begin{bullets}
  \item ローカルで動作する Overleaf クローン。Monaco エディタとライブ PDF プレビューの分割ビュー、ファイルツリー、ドラッグ＆ドロップアップロード、保存時自動コンパイルに対応。
\end{bullets}

\section{スキル}

\begin{tabular}{@{}p{0.28\linewidth}@{}p{0.72\linewidth}@{}}
  \textbf{プログラミング言語} & Python、OCaml、C、Java、TypeScript／JavaScript、Dart、SQL、Verilog \\
  \textbf{フレームワーク} & Next.js、React、FastAPI、Flutter、Express \\
  \textbf{データ・インフラ} & PostgreSQL、Firebase、GCP（Cloud Run、Cloud SQL）、Docker、Git \\
  \textbf{自動化} & Playwright、シェルスクリプト、Web スクレイピング \\
  \textbf{その他} & \LaTeX、Monaco、Linux、VS Code \\
\end{tabular}

\section{語学・資格}

\begin{tabular}{@{}p{0.28\linewidth}@{}p{0.72\linewidth}@{}}
  \textbf{語学} & 英語（流暢）・ビルマ語（母語）・日本語（JLPT N1）・スペイン語（日常会話）・中国語（基礎） \\
  \textbf{資格} & 日本語能力試験 N1（2023年8月取得） \\
\end{tabular}

\end{document}
